\documentclass[10pt]{article}
\usepackage[T1]{fontenc}
\usepackage[utf8]{inputenc}
\usepackage[french]{babel}
\usepackage{amsmath,amssymb,mathrsfs,mathpazo}
\usepackage{geometry}
\geometry{a4paper,margin=2cm}
\usepackage{graphicx,array,tabularx,tabularray,colortbl,multirow,multicol,booktabs,enumitem,pifont,ulem,fancybox,diagbox,cancel}
\usepackage{pstricks,pstricks-add,pst-plot,pst-tree,pst-node,pst-text,pst-func,pst-math,pst-bspline,pst-3dplot,pst-eucl,pst-solides3d}
\let\pscalculate\undefined
\usepackage{pst-bezier}
\usepackage{tikz,pgfplots,tkz-tab}
\usetikzlibrary{arrows,arrows.meta,patterns,calc,positioning,trees,decorations.markings,shapes.geometric,shapes.misc,shapes.arrows,intersections,decorations.pathmorphing,decorations.pathreplacing,backgrounds,shadows,babel}
\pgfplotsset{compat=1.18}
\usepgfplotslibrary{fillbetween}
\usepackage[np]{numprint}
\usepackage{listings,makecell,pas-tableur,twoopt}
\let\barre\undefined
\usepackage[pstricks]{tablvar}
\definecolor{darkgreen}{RGB}{0,100,0}
\providecommand{\red}{\color{red}}
\providecommand{\blue}{\color{blue}}
\providecommand{\green}{\color{green}}
\providecommand{\black}{\color{black}}
\providecommand{\white}{\color{white}}
\providecommand{\gray}{\color{gray}}
\providecommand{\cyan}{\color{cyan}}
\providecommand{\magenta}{\color{magenta}}
\providecommand{\yellow}{\color{yellow}}
\providecommand{\violet}{\color{violet}}
\providecommand{\orange}{\color{orange}}
\providecommand{\euro}{\texteuro}
\providecommand{\EUR}{\texteuro}
\providecommand{\ieme}{\textsuperscript{e}}
\providecommand{\degres}{\ensuremath{^\circ}}
\providecommand{\degree}{\ensuremath{^\circ}}
\providecommand{\vv}[1]{\overrightarrow{#1}}
\providecommand{\Coord}[1]{#1}
\providecommand{\decosix}{}
\providecommand{\starredbullet}{\ensuremath{\bullet}}
\pagestyle{empty}
\setlength{\parindent}{0pt}
\usepackage{adjustbox}\DeclareUnicodeCharacter{2236}{:}\DeclareUnicodeCharacter{2190}{\ensuremath{\leftarrow}}
\providecommand{\R}{\mathbb{R}}
\providecommand{\N}{\mathbb{N}}
\providecommand{\Z}{\mathbb{Z}}
\providecommand{\Q}{\mathbb{Q}}
\providecommand{\D}{\mathbb{D}}
\providecommand{\C}{\mathbb{C}}
\providecommand{\e}{{\mathrm{e}}}
\providecommand{\vect}[1]{\overrightarrow{#1}}
\providecommand{\vectt}[1]{\overrightarrow{\mathrm{#1}}}
\providecommand{\vvt}[1]{\overrightarrow{\mathrm{#1}}}
\providecommand{\np}[2][]{{#2~#1}}
\providecommand{\vv}[1]{{\overrightarrow{#1}}}
\renewcommand{\d}{{\mathrm{d}}}
\providecommand{\degres}{{^\circ}}
\providecommand{\degre}{{^\circ}}
\providecommand{\blue}{\color{blue}}
\providecommand{\red}{\color{red}}
\providecommand{\enskip}{\quad}
\providecommand{\llbracket}{\lbrack\!\lbrack}
\providecommand{\rrbracket}{\rbrack\!\rbrack}
\providecommand{\euro}{{\text{€}}}
\providecommand{\no}{{\text{n°}}}
\providecommand{\interval}[2]{{[#1;#2]}}
\providecommand{\intervFO}[2]{{[#1;#2[}}
\providecommand{\dotfill}{\ldots}

\providecommand{\ds}{\displaystyle}
\providecommand{\textding}[1]{#1}
\providecommand{\Oij}{\(\mathrm{O};\vec{\imath},\vec{\jmath}\)}
\providecommand{\Oijk}{\(\mathrm{O};\vec{\imath},\vec{\jmath},\vec{k}\)}
\providecommand{\Ouv}{\(\mathrm{O};\vec{u},\vec{v}\)}
\providecommand{\intervFF}[2]{\ensuremath{[#1;#2]}}
\renewcommand{\euro}{\ifmmode\text{\texteuro}\else\texteuro\fi}
\DeclareUnicodeCharacter{00B0}{\ifmmode^{\circ}\else\textdegree\fi}
\DeclareUnicodeCharacter{20AC}{\ifmmode\text{\texteuro}\else\texteuro\fi}
% Énoncés : Éducation nationale. Transcription et corrigés : APMEP, Denis Vergès et contributeurs. Découpage et classement : CoopMaths, Sébastien Lozano, Éric Elter et contributeurs. Adaptation de la présentation : Maths Explorer.
% CC BY-SA : https://coopmaths.fr/www/about/
% https://coopmaths.fr/alea/static/e3c/2020/tex/e3c_2020_00_sujet6_4_cor.tex
\begin{document}
\noindent\textit{APMEP / CoopMaths. CC BY-SA. Présentation adaptée par Maths Explorer.}\par\medskip
\subsection*{1. Exprimer le rayon de la base en fonction de $h$.}
	
	Le théorème de Pythagore appliqué au triangle rectangle d'hypoténuse la génératrice s'écrit :
	\[
	20^2 = h^2 + r^2 \iff r^2 = 400 - h^2 \iff r = \sqrt{400 - h^2}
	\]
	
	\subsection*{2. Démontrer que le volume du cône, en fonction de sa hauteur $h$, est : $V(h) = \dfrac{\pi}{3}(400h - h^3)$.}
	
	On a donc :
	\[
	V = \dfrac{1}{3} \pi (400 - h^2)h = \dfrac{\pi}{3}(400h - h^3).
	\]
	
	\subsection*{3. Quelle hauteur $h$ choisir pour que le volume du cône soit maximum ?}
	
	La fonction polynôme $V$ est dérivable et
	\[
	V'(h) = \dfrac{\pi}{3}(400 - 3h^2).
	\]
	
	Le signe de $V'(h)$ est celui de $400 - 3h^2$.
	\begin{itemize}
		\item $400 - 3h^2 > 0 \iff 400 > 3h^2 \iff h^2 < \dfrac{400}{3} \iff h < \sqrt{\dfrac{400}{3}}$.
		\item $400 - 3h^2 < 0 \iff 400 < 3h^2 \iff h^2 > \dfrac{400}{3} \iff h > \sqrt{\dfrac{400}{3}}$.
		\item $400 - 3h^2 = 0 \iff 400 = 3h^2 \iff h^2 = \dfrac{400}{3} \iff h = \sqrt{\dfrac{400}{3}}$.
	\end{itemize}
	
	Du signe de la dérivée résultent les variations de $V$ :
	\begin{itemize}
		\item Sur $[0; \sqrt{\dfrac{400}{3}}]$, la fonction est croissante.
		\item Sur $[\sqrt{\dfrac{400}{3}}; +\infty]$, la fonction est décroissante.
	\end{itemize}
	
	\[
	V\left(\sqrt{\dfrac{400}{3}}\right) = \dfrac{\pi}{3}\left(400 \times \sqrt{\dfrac{400}{3}} - \left(\sqrt{\dfrac{400}{3}}\right)^3\right) \approx \numprint{3224.53} \text{ cm}^3
	\]
	
	et ce pour une hauteur optimale de $\sqrt{\dfrac{400}{3}} \approx \numprint{11.55}$ cm.

\end{document}
